\documentclass[11pt]{article}
\usepackage{fontspec}
\usepackage{polyglossia}
\usepackage[utf8]{inputenc}
\usepackage{xunicode}
\setmainfont{Helvetica}
\usepackage[paperwidth=11in,paperheight=8.5in,margin=1in,headheight=0.0in,footskip=0.5in,includehead,includefoot,portrait]{geometry}
\usepackage[absolute]{textpos}
\TPGrid[0.5in, 0.25in]{23}{24}
\parindent=0pt
\parskip=12pt
\usepackage{nopageno}
\usepackage{graphicx}
\graphicspath{ {./images/} }
\usepackage{amsmath}
\usepackage{tikz}
\newcommand*\circled[1]{\tikz[baseline=(char.base)]{
            \node[shape=circle,draw,inner sep=1pt] (char) {#1};}}
            
\usepackage{xcolor}
\definecolor{mypink}{rgb}{0.572549019607843, 0.184313725490196, 0.203921568627451}
\definecolor{myred}{rgb}{0.843137254901961, 0.290196078431373, 0.270588235294118}
\definecolor{myblack}{rgb}{0.094117647058824, 0.223529411764706, 0.254901960784314}
\definecolor{mygrey}{rgb}{0.745098039215686, 0.725490196078431, 0.690196078431373}

\setmainlanguage{english}
\setotherlanguage{arabic}
\newfontfamily\arabicfont[Script=Arabic,Scale=0.75]{Amiri Bold}

\begin{document}

\begin{textblock}{23}(0, 1)
\begin{center}
\color{myblack}
\huge \textbf{FOREWORD}
\end{center}
\end{textblock}

\vspace*{0.25\baselineskip}

\begingroup

\begin{center}
\leftskip0.25in
The title of this work, \textbf{shqayq alnuemani}, is derived from the poem ``The Beloved Hemorrhaged Anemones'' by Mahmoud Darwish. In this poem, Darwish metaphorizes the red windflower of the Levant not only as a symbol for the heritage of land and peoplegroups but also deploys the image of a suddenly blooming hillside as a flood of blood, both in the sense of a cataclysmic wound and as a burst of new birth and future life. This musical composition attempts to sonically reconcile these tensions between life and death, nature and humanity, heritage, land, and violence. The musical textures frequently swing between beautiful, florid virtuosity and stark, unbending fauvist blocks. 
\rightskip\leftskip
\phantom{text} \hfill \phantom{()}

\vspace*{1.25\baselineskip}

\begin{textblock}{7.5}(7.5, 9)
\color{mypink}
\leftskip0.5in
The beloved hemorrhaged anemones.
The land of purple shimmered with his wounds.
The first of its songs: the blood of love that gods have shed;
and the last of them: blood...
\\
\[...\]
\\
The beloved hemorrhaged anemones.
The rocks of the hillside turned yellow
from the pain of difficult labor,
then reddened,
and the water flowed red
through the veins of our spring...
\\
\[...\]
\rightskip\leftskip
\\
\phantom{text} \hfill \phantom{()}
\\
(from `The Beloved Hemorrhaged Anemones'' by Mahmoud Darwish)

\end{textblock}


\end{center}

\endgroup

\vspace*{17\baselineskip}

\begin{center}
\color{myblack}
\huge \textbf{PERFORMANCE NOTES}
\end{center}

\color{black}
\leftskip0.25in
\textbf{Abbreviations} :accel. (accelerando), b.b. (behind bridge on wrapping), clb. (col legno battuto), clt. (col legno tratto), flaut. (flautando), gr. (gridato), N (normale), o.b. (on the wood of the bridge),  ord. (ordinario), pizz. (pizzicato), P (sul ponticello), MP (molto sul ponticello), XP (extreme sul ponticello),  rall. (rallantando), T (sul tasto), MT (molto sul tasto), and XT (extreme sul tasto).
\rightskip\leftskip
\phantom{text} \hfill \phantom{()}

\leftskip0.25in
\textbf{Accidentals} : After temporary accidentals, cancellation marks are printed also in the following measure (for notes in the same octave) and, in the same measure, for notes in other octaves. Accidentals are printed again if the same note appears later in the same measure – except if the note is immediately repeated
\rightskip\leftskip
\phantom{text} \hfill \phantom{()}

\leftskip0.25in
\textbf{Auxiliary Staves} : Some techniques in this work are notated by the use of auxiliary staves. The two examples within this piece are a staff with 4 lines and a staff with 9 lines. The 4-lined staff is deployed to represent the pseudo-barriolage bowing action of the right hand while the left hand (in a static fingering position) slides along the length of the strings. The 9-lined staff is a graphic representation of sul tasto to sul ponticello. Tasto is represented by dashed lines while ponticello is represented by solid lines. It is up to the performer's intuition whether the staff is read in a linear or logarithmic fashion. This 9-line staff is sometimes used to indicated complex arco bow trajectories and is sometimes used to indicate attack locations for col legno battuto.
\rightskip\leftskip
\phantom{text} \hfill \phantom{()}

\leftskip0.25in
\textbf{Tremolo} : Except when tremolo is notated above the staff (representing accelerando and ritardando), all tremoli should be played as fast as possible rather than as a measured subdivision.
\rightskip\leftskip
\phantom{text} \hfill \phantom{()}

\vspace*{2\baselineskip}

\begin{center}
\textit{shqayq alnuemani} was commissioned by VIVO Columbus.
\end{center}

\vspace*{2\baselineskip}

\begin{center}
\color{mypink}
duration: c. 9'30''
\end{center}

\end{document}
